% =============================================
% File: chapters/conclusion.tex
% KẾT LUẬN VÀ KIẾN NGHỊ — mục 3.9
% =============================================

\chapter*{KẾT LUẬN VÀ KIẾN NGHỊ}
\addcontentsline{toc}{chapter}{Kết luận và kiến nghị}
\label{chap:ketluan}


% --------------------------------------------------
\section*{1. Kết luận}
% --------------------------------------------------

Qua quá trình thực hiện đề tài \textit{``Nghiên cứu và ứng dụng bộ điều khiển PID
kết hợp mạng nơron nhân tạo (ANN) điều khiển Robot Bipedal Wheel sử dụng cơ cấu
5-bar parallel''}, nhóm nghiên cứu đã đạt được các kết quả sau:

\begin{enumerate}[leftmargin=1.5cm, label=\arabic*.]
    \item \textbf{Nhận dạng thành công hàm truyền động cơ NF5475} với độ chính xác
    95,36\% so với dữ liệu thực nghiệm, sử dụng công cụ \texttt{tfest} trong MATLAB.
    Hàm truyền bậc 4 xác định được phản ánh đúng đặc tính động học của động cơ
    ở nhiều dải tốc độ.

    \item \textbf{Thiết kế thành công bộ PI vận tốc} với tham số tối ưu tìm bằng
    giải thuật di truyền (GA). Bộ PI đạt độ vọt lố $< 5\%$ và thời gian xác lập
    $< 0{,}5$ s trên 9 mức setpoint từ 250 đến 2250 RPM.

    \item \textbf{Xây dựng và huấn luyện mạng ANN} dự đoán tham số $K_p$, $K_i$
    thích nghi theo sai số vận tốc. Mạng được xây dựng hoàn toàn từ đầu bằng Python,
    hội tụ sau 1000 epochs và không bị overfitting nhờ L2 regularization và Dropout.

    \item \textbf{Triển khai thực tế trên ESP32 WROOM-32} hệ thống điều khiển PI
    vận tốc với chu kỳ $\Delta t = 0{,}01$ s, hoạt động ổn định ở toàn dải setpoint
    thử nghiệm.

    \item \textbf{Phân tích động học cơ cấu 5-bar parallel}: xác định động học thuận,
    động học nghịch và ma trận Jacobian với thông số $A_1A_2 = 105$ mm,
    $L_{act} = 100$ mm, $L_{pas} = 180$ mm. Kết quả là nền tảng để mở rộng
    sang điều khiển quỹ đạo điểm cuối chân robot.
\end{enumerate}

Những kết quả trên chứng minh tính khả thi của việc kết hợp PID với ANN trong
điều khiển thích nghi động cơ DC trên nền tảng nhúng chi phí thấp, đồng thời
cung cấp quy trình nghiên cứu hoàn chỉnh và có thể tái hiện cho cộng đồng sinh viên.


% --------------------------------------------------
\section*{2. Hạn chế}
% --------------------------------------------------

Bên cạnh những kết quả đạt được, đề tài còn một số hạn chế cần thừa nhận:

\begin{enumerate}[leftmargin=1.5cm, label=\arabic*.]
    \item Mạng ANN chưa được kiểm chứng định lượng đầy đủ trên phần cứng thực;
    kết quả đánh giá hiện tại chủ yếu dựa trên dữ liệu offline.

    \item Thực nghiệm được thực hiện trong điều kiện tải cố định, chưa khảo sát
    hiệu năng khi tải thay đổi đột ngột hoặc có nhiễu từ môi trường ngoài.

    \item Bộ ANN chỉ dự đoán $K_p$, $K_i$ cho \textbf{một động cơ}; chưa mở rộng
    sang điều khiển đồng thời hai động cơ trong cơ cấu 5-bar parallel.

    \item Phần phân tích cơ cấu 5-bar chỉ dừng lại ở động học phẳng 2D, chưa
    bao gồm mô hình động lực học và chưa tích hợp thực tế vào robot hoàn chỉnh.

    \item Chưa có công bố khoa học chính thức từ kết quả đề tài.
\end{enumerate}


% --------------------------------------------------
\section*{3. Kiến nghị và hướng phát triển}
% --------------------------------------------------

Từ những kết quả và hạn chế trên, nhóm đề xuất các hướng phát triển tiếp theo:

\begin{enumerate}[leftmargin=1.5cm, label=\arabic*.]
    \item \textbf{Tích hợp ANN lên ESP32:} Nhúng mô hình ANN (sau khi lượng tử hóa
    hoặc xuất trọng số) trực tiếp lên ESP32, kiểm chứng thời gian suy luận
    và độ chính xác dự đoán trực tuyến (real-time).

    \item \textbf{Mở rộng sang điều khiển hai động cơ:} Áp dụng bộ PI--ANN cho
    cả hai động cơ trong cơ cấu 5-bar parallel, kết hợp với IK và Jacobian để
    điều khiển quỹ đạo điểm cuối chân robot.

    \item \textbf{Xây dựng robot bipedal wheel hoàn chỉnh:} Tích hợp hệ thống
    điều khiển vào phần cứng cơ khí thực, thực hiện thử nghiệm cân bằng và
    di chuyển trên bề mặt phẳng.

    \item \textbf{Nâng cao thuật toán điều khiển:} Khảo sát các phương pháp
    như Model Predictive Control (MPC), Reinforcement Learning hoặc
    Fuzzy--PID để so sánh và cải thiện hiệu năng bám setpoint.

    \item \textbf{Công bố kết quả:} Hoàn thiện số liệu thực nghiệm và chuẩn bị
    bài báo khoa học gửi đến tạp chí / hội nghị trong nước phù hợp.
\end{enumerate}
